\documentclass[11pt, a4paper, fontset=none]{ctexart}

\usepackage{assignment_style}

\begin{document}

% =========================================
% P288. Q1(1)
% =========================================
\begin{problem}{P288. Q1(1)}
    用对分区间法求方程
    \[
        x-\ln x=2
    \]
    在区间 $[2,4]$ 内的根，要求误差不超过 $10^{-5}$。
\end{problem}

\begin{solution}
    令
    \[
        f(x)=x-\ln x-2.
    \]
    因为
    \[
        f(2)=-\ln 2<0,\qquad
        f(4)=2-\ln 4>0,
    \]
    所以 $[2,4]$ 内存在根。

    对分法第 $n$ 次分半后的区间长度为
    \[
        \frac{b-a}{2^n}=\frac{2}{2^n}.
    \]
    为使区间长度不超过 $10^{-5}$，取
    \[
        \frac{2}{2^n}\le 10^{-5},
        \qquad n\ge \left\lceil \log_2(2\times 10^5)\right\rceil=18.
    \]
    经过 $18$ 次二分，得到
    \[
        x\approx 3.146191.
    \]
    因此所求根为
    \[
        \boxed{x\approx 3.146191}.
    \]
\end{solution}


% =========================================
% P288. Q2
% =========================================
\begin{problem}{P288. Q2}
    已知方程
    \[
        x^3-x^2-1=0
    \]
    在 $x_0=1.5$ 附近有根，试判断下列迭代格式：
    \[
        \text{(1)}\ x_{n+1}=1+\frac{1}{x_n^2},\qquad
        \text{(2)}\ x_{n+1}=\sqrt{\frac{1}{x_n-1}},
    \]
    \[
        \text{(3)}\ x_{n+1}=\sqrt[3]{1+x_n^2}
    \]
    在 $x_0=1.5$ 附近的收敛性。选一种迭代格式计算 $1.5$ 附近的根，要求结果具有 $5$ 位有效数字。
\end{problem}

\begin{solution}
    将三个迭代格式分别记为 $x=\varphi_i(x)$。

    \textbf{(1)}
    \[
        \varphi_1(x)=1+\frac{1}{x^2},\qquad
        \varphi_1'(x)=-\frac{2}{x^3}.
    \]
    在 $x=1.5$ 处，
    \[
        |\varphi_1'(1.5)|=\frac{2}{1.5^3}\approx 0.5926<1,
    \]
    故该迭代格式在根附近收敛。

    \textbf{(2)}
    \[
        \varphi_2(x)=(x-1)^{-1/2},\qquad
        \varphi_2'(x)=-\frac{1}{2}(x-1)^{-3/2}.
    \]
    在 $x=1.5$ 处，
    \[
        |\varphi_2'(1.5)|=\frac{1}{2(0.5)^{3/2}}\approx 1.4142>1,
    \]
    故该迭代格式在根附近不收敛。

    \textbf{(3)}
    \[
        \varphi_3(x)=(1+x^2)^{1/3},
    \]
    \[
        \varphi_3'(x)=\frac{2x}{3(1+x^2)^{2/3}}.
    \]
    在 $x=1.5$ 处，
    \[
        |\varphi_3'(1.5)|
        =\frac{3}{3(3.25)^{2/3}}
        \approx 0.4561<1,
    \]
    故该迭代格式在根附近收敛。

    选用第 (1) 个迭代格式，从 $x_0=1.5$ 出发，按
    \[
        x_{n+1}=1+\frac{1}{x_n^2}
    \]
    迭代，并以
    \[
        |x_{n+1}-x_n|<5\times 10^{-5}
    \]
    作为停止条件，得到
    \[
        x\approx 1.465556.
    \]
    因而该根的 $5$ 位有效数字为
    \[
        \boxed{x\approx 1.4656}.
    \]
\end{solution}


% =========================================
% P288. Q3(1)
% =========================================
\begin{problem}{P288. Q3(1)}
    用简单迭代法求方程
    \[
        x-\cos x=0
    \]
    的根，要求
    \[
        |x_{n+1}-x_n|<10^{-5}.
    \]
\end{problem}

\begin{solution}
    将方程写成等价的不动点形式
    \[
        x=\varphi(x)=\cos x.
    \]
    取区间 $[0,1]$，有
    \[
        f(0)=-1<0,\qquad f(1)=1-\cos 1>0,
    \]
    故根位于 $[0,1]$ 内。又
    \[
        \varphi([0,1])=[\cos 1,1]\subset [0,1],
    \]
    且
    \[
        |\varphi'(x)|=|\sin x|\le \sin 1<1,\qquad x\in[0,1].
    \]
    因此简单迭代在该区间内收敛。

    取
    \[
        x_0=0.5,\qquad x_{n+1}=\cos x_n.
    \]
    迭代至 $|x_{n+1}-x_n|<10^{-5}$ 时，得到
    \[
        x\approx 0.739082.
    \]
    所求根为
    \[
        \boxed{x\approx 0.739082}.
    \]
\end{solution}


% =========================================
% P288. Q7
% =========================================
\begin{problem}{P288. Q7}
    设有解方程
    \[
        12-3x+2\cos x=0
    \]
    的迭代法
    \[
        x_{n+1}=4+\frac{2}{3}\cos x_n.
    \]
    \begin{enumerate}[(1)]
        \item 证明对 $\forall x_0\in\R$，此迭代法均收敛到方程的根；
        \item 取 $x_0=4$，用此迭代法求方程根的近似值，使其误差不超过 $\dfrac{1}{2}\times 10^{-5}$。
    \end{enumerate}
\end{problem}

\begin{solution}
    记
    \[
        \varphi(x)=4+\frac{2}{3}\cos x.
    \]
    对任意 $x\in\R$，有
    \[
        \varphi(x)\in \left[4-\frac{2}{3},4+\frac{2}{3}\right]
        =\left[\frac{10}{3},\frac{14}{3}\right]\subset \R.
    \]
    又
    \[
        \varphi'(x)=-\frac{2}{3}\sin x,
    \]
    因而
    \[
        |\varphi'(x)|\le \frac{2}{3}<1,\qquad x\in\R.
    \]
    由压缩映射原理，迭代
    \[
        x_{n+1}=\varphi(x_n)
    \]
    对任意初值 $x_0\in\R$ 均收敛到唯一不动点 $x^\ast$。该不动点满足
    \[
        x^\ast=4+\frac{2}{3}\cos x^\ast,
    \]
    即
    \[
        12-3x^\ast+2\cos x^\ast=0.
    \]
    故此迭代法对任意初值均收敛到方程的根。

    取 $x_0=4$ 进行迭代。由上式可取压缩常数
    \[
        q=\frac{2}{3}.
    \]
    后验误差估计为
    \[
        |x^\ast-x_{n+1}|
        \le \frac{q}{1-q}|x_{n+1}-x_n|
        =2|x_{n+1}-x_n|.
    \]
    为使误差不超过 $\frac{1}{2}\times 10^{-5}$，只需令
    \[
        |x_{n+1}-x_n|<2.5\times 10^{-6}.
    \]
    迭代计算得
    \[
        x\approx 3.347403.
    \]
    因而所求近似根为
    \[
        \boxed{x\approx 3.347403}.
    \]
\end{solution}


% =========================================
% P288. Q10(1)
% =========================================
\begin{problem}{P288. Q10(1)}
    用 Newton 法求解方程
    \[
        x-\ln x=2,
    \]
    要求
    \[
        |x_{n+1}-x_n|<10^{-6}.
    \]
\end{problem}

\begin{solution}
    令
    \[
        f(x)=x-\ln x-2,\qquad f'(x)=1-\frac{1}{x}.
    \]
    Newton 迭代公式为
    \[
        x_{n+1}
        =x_n-\frac{f(x_n)}{f'(x_n)}
        =x_n-\frac{x_n-\ln x_n-2}{1-\frac{1}{x_n}}.
    \]
    由于
    \[
        f\left(\frac18\right)>0,\qquad f(1)<0,
    \]
    可取
    \[
        x_0=\frac18.
    \]
    迭代结果如下：
    \[
        \begin{array}{c|cc}
            n & x_n & f(x_n) \\ \hline
            0 & 0.125000 & 2.0444\times 10^{-1} \\
            1 & 0.154206 & 2.3672\times 10^{-2} \\
            2 & 0.158522 & 3.8451\times 10^{-4} \\
            3 & 0.158594 & 1.0437\times 10^{-7} \\
            4 & 0.158594 & 7.5495\times 10^{-15}
        \end{array}
    \]
    且
    \[
        |x_4-x_3|=1.9672\times 10^{-8}<10^{-6}.
    \]
    故所求根为
    \[
        \boxed{x\approx 0.158594}.
    \]
\end{solution}


% =========================================
% P288. Q13(1)
% =========================================
\begin{problem}{P288. Q13(1)}
    若 $x^\ast$ 是方程 $f(x)=0$ 的 $r$ 重根，即
    \[
        f(x)=(x-x^\ast)^r\varphi(x),\qquad \varphi(x^\ast)\ne 0.
    \]
    证明：Newton 法在 $x^\ast$ 邻近是线性收敛的。
\end{problem}

\begin{myproof}
    设
    \[
        e_n=x_n-x^\ast.
    \]
    因为
    \[
        f(x)= (x-x^\ast)^r\varphi(x),
    \]
    所以
    \[
        f'(x)
        =(x-x^\ast)^{r-1}
        \left[r\varphi(x)+(x-x^\ast)\varphi'(x)\right].
    \]
    Newton 法为
    \[
        x_{n+1}=x_n-\frac{f(x_n)}{f'(x_n)}.
    \]
    于是
    \[
        e_{n+1}
        =e_n-\frac{e_n^r\varphi(x_n)}
        {e_n^{r-1}\left[r\varphi(x_n)+e_n\varphi'(x_n)\right]}.
    \]
    整理得
    \[
        e_{n+1}
        =e_n
        \frac{(r-1)\varphi(x_n)+e_n\varphi'(x_n)}
        {r\varphi(x_n)+e_n\varphi'(x_n)}.
    \]
    当 $x_n\to x^\ast$ 时，$e_n\to 0$，且 $\varphi(x^\ast)\ne 0$，故
    \[
        \lim_{n\to\infty}\frac{|e_{n+1}|}{|e_n|}
        =\frac{r-1}{r}.
    \]
    对 $r>1$，有
    \[
        0<\frac{r-1}{r}<1,
    \]
    因而 Newton 法在重根 $x^\ast$ 邻近为线性收敛，收敛因子为
    \[
        \boxed{\frac{r-1}{r}}.
    \]
\end{myproof}


% =========================================
% P288. Q16(1)
% =========================================
\begin{problem}{P288. Q16(1)}
    试用割线法求解方程
    \[
        x-\ln x=2,
    \]
    要求
    \[
        |x_{n+1}-x_n|<10^{-5}.
    \]
\end{problem}

\begin{solution}
    令
    \[
        f(x)=x-\ln x-2.
    \]
    割线法迭代公式为
    \[
        x_{n+1}
        =x_n-\frac{f(x_n)(x_n-x_{n-1})}{f(x_n)-f(x_{n-1})}.
    \]
    取实数定义域内的初值
    \[
        x_0=\frac18,\qquad x_1=\frac14.
    \]
    此时
    \[
        f\left(\frac18\right)>0,\qquad f\left(\frac14\right)<0,
    \]
    两初值夹住小根。

    迭代结果如下：
    \[
        \begin{array}{c|cc}
            n & x_n & f(x_n) \\ \hline
            0 & 0.125000 & 2.0444\times 10^{-1} \\
            1 & 0.250000 & -3.6371\times 10^{-1} \\
            2 & 0.169980 & -5.7945\times 10^{-2} \\
            3 & 0.154815 & 2.0338\times 10^{-2} \\
            4 & 0.158755 & -8.5212\times 10^{-4} \\
            5 & 0.158597 & -1.2094\times 10^{-5} \\
            6 & 0.158594 & 7.2820\times 10^{-9}
        \end{array}
    \]
    且
    \[
        |x_6-x_5|=2.2809\times 10^{-6}<10^{-5}.
    \]
    故所求根为
    \[
        \boxed{x\approx 0.158594}.
    \]
\end{solution}

\end{document}
